% 张三慧《大学物理学（第4版）》力学、热学，第2章运动与力
% 原书章末习题2.1--2.12；教材PDF物理页82--84（印刷页66--68）。
% 图2.32--2.40保留准确几何关系的文字说明；原图尚需矢量重绘。
% 单章候选源，不得进入bank/正式高难综合题库或自动组卷。
\section*{第2章\quad 运动与力}
\subsection*{习题（2.1--2.12；后续题号待录入）}

\paragraph{习题 2.1}
用力 \(F\) 推水平地面上一质量为 \(M\) 的木箱（原书图2.32）。
设力 \(F\) 与水平面的夹角为 \(\theta\)，
木箱与地面间的滑动摩擦系数和静摩擦系数分别为
\(\mu_k\) 和 \(\mu_s\)。
\begin{enumerate}
\item 要推动木箱，\(F\) 至少应多大？此后维持木箱匀速前进，
\(F\) 应需多大？
\item 证明当 \(\theta\) 大于某一值时，无论用多大的力
\(F\) 也不能推动木箱。此 \(\theta\) 角是多大？
\end{enumerate}
\textit{图2.32：}木箱在水平地面上，外力指向右下，
与向右的水平线夹角 \(\theta\)。

\paragraph{习题 2.2}
设质量 \(m=0.50\,\mathrm{kg}\) 的小球挂在倾角 \(\theta\)
的光滑斜面上（原书图2.33）。
\begin{enumerate}
\item 当斜面以加速度 \(a=2.0\,\mathrm{m/s^2}\)
沿如图2.33所示的方向运动时，绳中的张力及小球对斜面的正压力各是多少？
\item 当斜面的加速度至少为多大时，小球将脱离斜面？
\end{enumerate}
\textit{图2.33：}斜面向右上方倾斜，绳沿斜面系在小球上方，
斜面整体向右加速；原题未给定 \(\theta\) 的具体数值，
答案保留 \(\theta\) 参数。

\paragraph{习题 2.3}
一架质量为 \(5000\,\mathrm{kg}\) 的直升机吊起一辆
\(1500\,\mathrm{kg}\) 的汽车以 \(0.60\,\mathrm{m/s^2}\)
的加速度向上升起。
\begin{enumerate}
\item 空气作用在螺旋桨上的上举力多大？
\item 吊汽车的绳绳中张力多大？
\end{enumerate}

\paragraph{习题 2.4}
如原书图2.34所示，一个高楼擦窗工人利用一定滑轮自己控制下降。
\begin{enumerate}
\item 要自己慢慢匀速下降，他需要用多大力拉绳？
\item 如果他放松些，使拉力减少 \(10\%\)，
他下降的加速度将多大？设人和吊桶的总质量为 \(75\,\mathrm{kg}\)。
\end{enumerate}
\textit{图2.34：}定滑轮一端连接吊桶，另一端由吊桶中工人向下拉住；
忽略滑轮、绳的质量和摩擦。

\paragraph{习题 2.5}
原书图2.35中 \(A\) 为定滑轮，\(B\) 为动滑轮，
3个物体的质量分别为
\(m_1=200\,\mathrm g,\ m_2=100\,\mathrm g,\
m_3=50\,\mathrm g\)。
\begin{enumerate}
\item 求每个物体的加速度。
\item 求两根绳中的张力 \(T_1\) 和 \(T_2\)。
假定滑轮和绳的质量以及细绳的伸长和摩擦力均可忽略。
\end{enumerate}
\textit{图2.35：}第一根绳跨过上方定滑轮 \(A\)，
左端挂 \(m_1\)，右端悬挂下方动滑轮 \(B\)；
第二根绳跨过 \(B\)，两端分别挂 \(m_2,m_3\)。

\paragraph{习题 2.6}
在一水平的直路上，一辆车速为
\(v=90\,\mathrm{km/h}\) 的汽车的刹车距离为
\(s=35\,\mathrm m\)。
如果路面相同，只是有 \(1:10\) 的下降斜度，
这辆汽车的刹车距离将变为多少？

\paragraph{习题 2.7}
桌上有一质量 \(M=1.50\,\mathrm{kg}\) 的板，
板上放一质量 \(m=2.45\,\mathrm{kg}\) 的另一物体。
物体与板、板与桌面之间的摩擦系数均为 \(\mu=0.25\)。
要将板从物体下面抽出，至少需要多大的水平力？

\paragraph{习题 2.8}
如原书图2.36所示，在一质量为 \(M\) 的小车上放一质量为
\(m_1\) 的物块，它用细绳通过固定在小车上的滑轮与质量为
\(m_2\) 的物块相连；物块 \(m_2\) 靠在小车的前壁上
而使悬线竖直，忽略所有摩擦。
\begin{enumerate}
\item 当用水平力 \(F\) 推小车使之沿水平桌面加速前进时，
小车的加速度多大？
\item 如果要保持 \(m_2\) 的高度不变，力 \(F\) 应多大？
\end{enumerate}
\textit{图2.36：}小车被水平推向右，\(m_1\) 在车顶面，
绳经车右上端的滑轮后竖直悬挂贴着车前壁的 \(m_2\)。

\paragraph{习题 2.9}
按照距离约 \(38\times10^5\,\mathrm{km}\) 外的地球上
飞行控制中心发来的指令，点燃自身的制动发动机后，
我国第一颗月球卫星“嫦娥一号”于2007年11月7日正式进入科学探测工作轨道
（原书图2.37）。该轨道为圆形，离月面的高度为
\(200\,\mathrm{km}\)。求“嫦娥一号”的运行速率（相对月球）
与运行周期。
\textit{需结合本书数值表中的月球半径和质量计算。}

\paragraph{习题 2.10}
两根弹簧的劲度系数分别为 \(k_1\) 和 \(k_2\)。
\begin{enumerate}
\item 试证明它们串联起来时（原书图2.38(a)），
总的劲度系数为
\[
 k=\frac{k_1k_2}{k_1+k_2}.
\]
\item 试证明它们并联起来时（原书图2.38(b)），
总的劲度系数为
\[
 k=k_1+k_2.
\]
\end{enumerate}

\paragraph{习题 2.11}
如原书图2.39所示，质量
\(m=1200\,\mathrm{kg}\) 的汽车在一弯道上行驶，
速率 \(v=25\,\mathrm{m/s}\)。
弯道的水平半径 \(R=400\,\mathrm{m}\)，
路面外高内低，倾角 \(\theta=6^\circ\)。
\begin{enumerate}
\item 求作用于汽车上的水平法向力与摩擦力。
\item 如果汽车轮与轨道之间的静摩擦系数
\(\mu_s=0.9\)，要保证汽车无侧向滑动，
汽车在此弯道上行驶的最大允许速率应是多大？
\end{enumerate}

\paragraph{习题 2.12}
现已知木星有16个卫星，其中4个较大的是伽利略用他自制的望远镜
在1610年发现的（原书图2.40）。
这4个“伽利略卫星”中最大的是木卫三，
它到木星的平均距离是 \(1.07\times10^6\,\mathrm{km}\)，
绕木星运行的周期是 \(7.16\,\mathrm d\)。
试由此求出木星的质量，忽略其他卫星的影响。
